\documentclass[12pt,a4paper]{book}

\usepackage[utf8]{inputenc}
\usepackage[T1]{fontenc}
\usepackage{geometry}
\geometry{margin=1in}

\usepackage{titlesec}
\usepackage{parskip}
\usepackage{enumitem}

\usepackage{xcolor}
\usepackage{tcolorbox}
\usepackage{fancyhdr}
\tcbuselibrary{skins}

\usepackage{epigraph}


% Chapter style
\titleformat{\chapter}[display]
  {\normalfont\Large\bfseries\centering}
  {\chaptertitlename\ \thechapter}{1em}{}
\titlespacing*{\chapter}{0pt}{50pt}{40pt}

% Section / Subsection style
\titleformat{\section}
  {\normalfont\large\bfseries}
  {\thesection}{1em}{}
\titleformat{\subsection}
  {\normalfont\normalsize\bfseries}
  {\thesubsection}{1em}{}

% Remove "Chapter X" from header
\pagestyle{plain}
\definecolor{forgecolor}{RGB}{192, 192, 192}   % Light Gray (recommended)

\newtcolorbox{scripturetopbox}{
  colback=forgecolor!8,
  colframe=forgecolor,
  boxrule=1.2pt,
  sharp corners,
  enhanced,
  width=\textwidth,
  top=4pt,
  bottom=4pt,
  valign=center,
  halign=left,
  fontupper=\bfseries\small,
  nobeforeafter
}

% Custom command — call this at the top of any page/chapter you want the box
\newcommand{\scripturetop}[1]{
  \begin{scripturetopbox}
  #1
  \end{scripturetopbox}
  \vspace{6pt}   % small space after the box
}

\usepackage{needspace}
\clubpenalty=10000
\widowpenalty=10000


\begin{document}

\thispagestyle{fancy}   % ← Add this line so the very first page also gets the box

% Title page
\begin{titlepage}
\begin{center}
\vspace*{3cm}   % reduced from 3cm

{\Huge \bfseries Shadow to Steel}

\vspace{0.8cm}

{\Large Obedience, Iron, and the Return of the King}

\vspace{1.2cm}

{\Large Glenn Michael}

\vspace{1.8cm}

\epigraph{The King does not return the same as when he left. 
He returns forged in the fire of obedience.}{}

\vspace{1.5cm}
\end{center}
\end{titlepage}

% Dedication (optional)
\thispagestyle{empty}
\begin{center}
  \vspace*{8cm}
  {\itshape
    To those walking through the valley and the forge,\\
    may you walk the path that leads to steel.
  }
\end{center}
\clearpage

% Table of Contents
\tableofcontents
\clearpage

% Introduction / Preface (optional)
\chapter*{Introduction}

%\addcontentsline{toc}{chapter}{Introduction}

% Custom command — call this at the top of any page/chapter you want the box
\scripturetop{
If you make a vow to the Lord your God, do not be lax in paying it off, for the Lord your God will certainly require it of you and you would be guilty of sin. - Deuteronomy 23:21
}

Following years of decline, on July 5, 2022 I was diagnosed with Relapsing Polychondritis — a rare, chronic, and often fatal autoimmune disease that attacks and destroys the body’s cartilage, particularly in the eyes, ears, nose, respiratory system, and trachea.

In my case, the disease had ravaged nearly every system. It destroyed my right knee (requiring surgery), my left hip (requiring replacement), and my nasal cartilage (rebuilt from chest cartilage, which revealed almost none left in the ribs). Looking back now, that outcome was merciful compared to the total destruction and floating chest many others endure.

In the recliner,  facing what looked like my certain end, a great fear fell over me: I was going to die. After a few minutes, I thought to myself, “This changes nothing. I was born,  and like every person born into this world,  I am going to die.” The only difference now was that I had an idea of the timing — imminent — and the means: destruction of my cartilage and airways.

Then I prayed to the Lord. First, I confessed my sins against Him. Second, I apologized for my arrogance in believing all I had accomplished was due to me alone, instead of recognizing those things as blessings from the Lord. Finally, I said, “If it is my time, then it is my time. If it is not my time, tell me what to do to fix myself, and I’ll do it.”

The next morning, in the twilight between sleep and waking, the Lord answered: “Bench. Squat. Deadlift.”

For those unfamiliar with powerlifting or gym culture, these are known as the Big Three — multi-joint compound lifts. A counter-intuitive recommendation for a patient with a serious connective tissue disease.

That answer was not a miracle. It was a path. A forge. A way through the valley of the shadow of death.  This book is the story of that path — from near-death weakness to restoration,  from betrayal to sovereignty,  from Shadow to Steel.

% Chapter 1 – The Dream
\chapter{The Black Marble Room}
Five years before the recliner, before the diagnosis, before the voice, I was already in the Valley of the Shadow of Death. I just didn’t know the name for it yet.

Nothing was going right.

\begin{itemize}

\item I had been wrongfully terminated from my job, my restricted stock awards, spitefully canceled.  I lost a significant amount of wealth.

\item Shortly afterward, lightning struck our home. The top floor burned. The lower floor flooded from the firefighting efforts. Most of the house was destroyed.

\item Having lost my job, the person that had driven the termination was using the full legal weight of the bank which I had worked to blacklist me, I began an independent consultancy which required me to be away from home five to six days a week,

\item My family relationships began to fray under the strain.

\end{itemize}

Living alone in cheap hotel rooms sinking with neither an end nor aid in sight the weight of it all became unbearable.  One night, I cried out to the Lord with raw desperation: “Lord,  I cannot bear this any longer.   If there is an end to it, please let me know.   If not, at least tell me that too, so I can resign myself to this trap.”

That night I had the most vivid dream of my life.

I stood in a vast,  windowless chamber carved from a single block of polished black marble.  The walls and floor were seamless, cold, reflective - like night itself made solid.  There was no door, no window, no visible exit. Far above, impossibly high in the ceiling, a tiny opening let in a faint shaft of light.  It might as well have been on the moon.

I tried to climb.  I wedged my body into the smooth corners, straining upward, but the marble offered no grip. Even if I gained height, exhaustion would have pulled me back down to my death.

Defeated and breathing hard, a simple marble bench appeared behind me. I sat down and told myself, “Don’t panic.  You’re a smart person.  You can figure this out.”

Then I felt a presence to my left — masculine, calm, and authoritative — yet I could never see Him directly.  He remained just beyond the edge of my vision, always slightly out of sight.  In a quiet, steady voice He said, “We have decided to show you what is going to happen to you.”  

He led me to the left-most corner of the room — the very corner I had already examined and found nothing. Now there was a doorway. Behind it stretched a long hallway lined with identical black marble doors on both sides. We walked until He stopped at one and said, “This is your door.”

I opened it and stepped into a smaller black marble room, identical to the first. In the center, on a raised pedestal, sat a golden tabernacle — exactly like the one that holds the Eucharist in church. I opened the doors and looked inside.

I cannot describe what I saw. I can only tell you the single overwhelming thought that flooded me: That’s not so bad.

I stepped back into the hallway. The presence was still there. We walked to the far end and exited. I was back in the original large chamber. The marble bench remained, but the doorway and hallway had vanished. Only the tiny unreachable light high above remained.

I sat down again, still trying to solve the puzzle, and then I woke up.

I never forgot that dream. It was the first quiet answer to a prayer I had not yet fully prayed. Seven years later, when I sat in the recliner facing what looked like certain death, the Lord had already shown me the truth: even in the most sealed, lightless chamber, a way out existed — if I was willing to wait, to listen, and to obey.

The dream was the seed planted in darkness.
The prayer — “tell me what to do and I will do it” — became the vow.

And the iron would become the path.

I would come to realize, that I was walking through the valley of the shadow of death.  But, before that realization came to me another was to precede.  The same friend that had asked if I have made a bargain with or a vow to the Lord had also likened my circumstance to that of Job.  

The Book of Job is one of the oldest and most profound books in the Bible.  It is not primarily about why "bad things happen to good people."  It is about \textbf{faith under extreme testing} - whether a man will serve God for who God is, or only for the blessings he receives.  The story opens in Heaven with a dialog between God and Satan.  God points out Job as "blameless and upright, a man who fears God and shuns evil."  Satan immediately accuses Job of conditional faith: "Does Job fear God for nothing?" That is to say, Satan claims Job only serves God because God has built a protective hedge of blessing around Job.

God then permits Satan to test Job - but with strict limits.  Satan cannot kill Job or take his life.  Everything else is on the table.  The story illustrates that Satan does not attack randomly.  Rather, Satan uses a systematic, multi-layered strategy designed to break Job's faith from every direction.  And rest assured Satan will use the same strategy with each and every person.  Satan's attack vectors as outlined in chapters one through two are outlined below:
\newpage
\begin{enumerate}
\item \textbf{Economic - Livelihood Destruction (Job 1:14-17)}
\begin{itemize}
\item Sabeans steal oxen and donkeys and kill servants.
\item Fire from heaven burns the sheep and shepherds.
\item Chaldeans steal camels and kill more servants.
\end{itemize}
\textit{Satan wipes out Job's entire wealth and business in one day.  The goal: remove security and provision.}
\item \textbf{Family - Emotional Devastation (Job 1:18-19}
\begin{itemize}
\item A great wind collapses the house where all ten of Job's children are feasting and kills everyone of them.
\end{itemize}
\textit{Satan removes Job's entire legacy and support system in a single moment.  The goal crush the heart.}
\item \textbf{Physical - Bodily Affliction (Job 2:7-8)}
\begin{itemize}
\item Satan strikes Job with painful sores from the soles of his feet to the top of his head. 
\end{itemize}
\textit{Job is left scraping himself with broken pottery while sitting in ashes.  The goal: constant, unrelenting torment that makes every moment miserable}
\item \textbf{Relations - Social Isolation (Job 2:9-10 and throughout)}
\begin{itemize}
\item Job's wife tells him to "Curse God and die."
\item Job's three friends (and later Elihu) arrive, but instead of comfort they accuse Job of hidden sin and judge him relentlessly for 30+ chapters.
\end{itemize}
\textit{Satan turns the people closest to Job into sources of pain and accusation.  The goal total loneliness.}
\item \textbf{Spiritual - Psychological Attack (the entire of the book)}
\begin{itemize}
\item Job is plunged into deep despair, questioning why he was ever born, feeling abandoned by God, and wrestling with the apparent injustice of it all.
\end{itemize}
\textit{Satan attacks the mind and the spirit, trying to drive Job to curse God.  The goal: break the relationship with God}

\end{enumerate}

I had no understanding at the time that I was living through the very same multi-layered attack Satan had used against Job.  In my ignorance, I simply fought harder to keep my head above water. The more I struggled, the deeper I sank into a darkness I could not name or escape.
I was already walking through the valley of the shadow of death.

What I did not yet realize was that the Valley was also the forge.


\chapter{The Valley – The Place of Shadow and Fear}
\scripturetop{
2 My God, my God, why have you forsaken me? \\
\hspace*{10pt} Why have you paid no heed to my cry for help, to my cries of anguish?\\ 
3 O my God, I cry by day, but you do not answer,by night,  \\
\hspace*{10pt} but I am afforded no relief. \\
20 But you, O lord, do not remain aloof from me;\\
\hspace*{10pt} O my Strength, come quickly to my aid. -- Psalm 22
}

Psalm 22 was written by King David around 1000–970 BC during one of his darkest periods. Most scholars believe it was composed while he was on the run during the rebellion of his own son Absalom. At the time, David was hunted, betrayed by his closest advisor, publicly humiliated, and utterly alone — living in the valley of the shadow of death.

I knew exactly how he felt.

For years the misfortune had been grinding me down. I felt a complete emotional disconnect from the woman I loved. The marriage that once anchored me had gone cold and silent. I prayed the same desperate question David did: My God, why have you abandoned me? At first I convinced myself the Lord had left because I was unworthy—I hadn’t been to church in years. Then I began to wonder if God existed at all. Neither thought brought comfort.

Standing alone on the train platform into town, the sky darkened by heavy clouds. It was cold and damp, and a stinging bitter wind whipped across my face. At that moment I had an epiphany: this is exactly how my inner world felt. Yet I continued to pray. I still believed in God. I still sought comfort from the Lord. The devastating truth was not that God did not exist, but that I could no longer feel His presence. And I had no idea that I would have to walk through the valley of the shadow of death before I would feel it again.

The Summer of 2021, I was walking through the valley- I didn't fully recognize it.  June.  Middle of the season, yet I wore a hoodie outdoors because I was freezing down to my core. My right knee had been painful and swollen for six months due to a torn meniscus. - how it happened was beyond me.  After repeated cortisone injections, insurance finally approved surgery.   During the follow-up the surgeon told me he had sent fluid from my knee for cancer testing—it looked that abnormal. The result came back negative.   Only later visit would he realize the fluid was classic for the autoimmune disease that was already destroying me.

Physical therapy after knee surgery brought new pain in my left hip.   Within weeks I could barely walk.  The hip surgeon showed me comparative X-rays: the cartilage in my left hip had been completely erased while the right hip remained unchanged. “What could do something like that?” I asked. He had no answer.

Then came the dry cough and shortness of breath—constant, unnatural.   Interns and nurses noticed it.  The hip replacement itself triggered a major flare.   Nasal cartilage was destroyed, leaving me with a saddle-nose deformity.   The shortness of breath and chest pain continued.  I bought a Peloton, mostly to help build leg strength following the two surgeries and I believe the shortness of breath convinced the was due to poor conditioning - little did I know that the kind of shortness of breath I was experiencing was a classic autoimmune symptom. Every ride became a battle: coughing, gasping, fighting an invisible enemy that was killing me slowly, methodically, day by day.

January 2022, late one night I laying in bed facing the truth: a silent killer was winning.   I turned to the Lord with a simple, stripped-down prayer: “Lord, something is killing me.   All I want is to know - what it is.   Please get me in front of someone who can tell me before I die.”

The next morning, in that quiet twilight between sleep and waking, a voice—not my own—spoke with calm authority:

“Don’t worry about work, don’t worry about insurance, just call everyone you can and take whatever appointment is available—today, tomorrow, a month, a year from now - It doesn’t matter.”

I started with the obvious: my ruined nose. I called ENT offices until, on the fifth try, a cheerful woman at Rush University answered. “You’re in luck—we just reopened this week.” She booked me with the department’s top specialist, Dr. Papagiannopoulos.

A week later, I was sitting in the exam room waiting for the doctor.  Once he arrived and we traded pleasantries, I went over the litany of maladies that had befallen me, each in its proper order.  I dropped the surgical mask they had given me to wear at admittance showing him the saddle nose deformity and said - "then this happened."  He replied, "I am sorry that happened to you."  This was the first time anyone had expressed sympathy for the deformity that had befallen me.  after which, he immediately referred me to rheumatology.  

That same day, I met Dr. Block.  He entered the exam room, looked me over and asked who referred me.  I replied, Dr. Papagiannopoulos, he nodded in approval.  I had managed to find the right doctor, one that could see the seriousness of my situation, and who would in turn refer me rightly to another that would help me - my prayer had been answered.

Dr. Block agreed, we needed to repair my nose.  A biopsy was needed to aid in the identification of the disease.  Over the course of the next six months, while waiting for my nasal reconstruction, I returned to Rush on a monthly basis for routine examinations and blood draws.   During this time I met Doctor Revenaugh the surgeon that would repair my nose.  Initially, he indicated that he would repair the saddle deformity but not the hole in the septum.   I explained to him, I realized the septum repair had at best and 80 percent success rate but even in the case of a failure it is an improvement.  I asked him to try, if it does not fully succeed that is fine at least we tried.  He agreed to try.

The surgery went forward, I splints in my nose to protect the repair work to my septum, they would remain for the summer.  The chest cartilage recovered to repair the saddle deformity was, as Dr. Revenaugh put it - 'barely viable'.  The disease had destroyed my chest cartilage.  Fortunately, it had remodeled itself into a fibrous tissue, something they had never seen before; I had been spared the fate of a 'loose chest', which is common in these instances.  

On my next visit in July they delivered the diagnosis Relapsing Polychondritis; until the surgery the disease had been running at a constant low level - insidious - making it unidentifiable.  It had impacted the cartilage in my right knee, destroyed the cartilage in my nose, chest, left hip, and was now slowly destroying the cartilage in my respiratory tract and trachea.  

My childhood asthma returned and breathing became increasingly difficult.  Doctor Block referred me for a pulmonary test.  At the follow up, I met doctor Vashi, she began treating me for both my resurgent asthma as well as the damage done to my respiratory tract.

The diagnosis landed like a sentence: Relapsing Polychondritis. I now knew the name of the silent killer that had been destroying me piece by piece.  My earlier prayer had been answered.  In the recliner that July evening, the weight of it settled fully.  I was going to die.

After a few minutes the fear gave way to clarity. I had already walked through the valley. I had already seen the black marble room. I had already heard the first quiet answer. In that moment I prayed as David once did in Psalm 22:

I confessed my sins against the Lord. I apologized for the arrogance of believing everything I had accomplished was by my own strength instead of recognizing those things as blessings from Him. Then I made no bargain — only a vow.

“If it is my time, then it is my time.  If it is not my time, tell me what to do to fix myself, and I will do it.”

The next morning, in the twilight between sleep and waking, the voice returned — calm and unmistakable:

“Bench.  Squat.  Deadlift.”

For those unfamiliar with powerlifting, these are the Big Three — the foundational compound lifts. For a man whose cartilage was being systematically destroyed by disease, the prescription sounded absurd. Yet in that moment, I understood: this was not a miracle. This was a path. This was the forge — the forge I knew so well from childhood.

The Valley had shown me the darkness. The Lord had shown me the way out.  Now the iron would become the fire that would refine the shadow of death into steel.

Even though I walk through the valley of the shadow of death, I will fear no evil, for you are with me; your rod and your staff, they comfort me. — Psalm 23:4-5

The Lord heard my cries and answered — not with a miracle, but with a prescription.  I was to step from the valley of the shadow of death into the forge that refines, and I knew the Lord was with me.

\chapter{The Forge – The Fire That Refines}

\scripturetop{
2 But who will be able to endure the day of his coming, \newline
\hspace *{10pt} and who can stand when he appears?\newline
For he is like a refiner's fire \newline
\hspace*{10pt} or fuller's soap.\newline
He will sit refining and purifying:\newline
\hspace*{10pt} he will purify the descendants of Levi\newline
and refine them like gold and silver.\newline
\hspace*{10pt} so that they may in righteousness\newline
\hspace*{10pt} offer due sacrifice to the Lord.  -- Malachi 3:2-3\\
}

I have known the fire of the forge longer than I realized.  As a child, I suffered from severe asthma. Each day was the same: some days were manageable, others a struggle.  But the nights were different.  I would wake in the late hours gasping for air, chest tight, lungs constricted.  My mother would come to my aid.  Some nights I would pass out from lack of oxygen. Other nights I would collapse from exhaustion simply trying to breathe.

I never complained.  I just kept going.

Years later, my Uncle Mike — who suffers from COPD — reflected on those days after Sunday brunch.  He looked at me and said, “My sister once told me, you never complained about your asthma. You just kept going.”

He was right.  Even then, without knowing the neither name nor its purpose, I was already walking through the first forge.  The same enemy that would later attack my cartilage,  my marriage,  and my peace was already testing me in the quiet hours of childhood.  He could not break me then, and he would not break me later. 
 
As a child, I would often pray to the Lord, why this affliction, why me?  Always, to no answer, not because the Lord did not hear my prayer but because the prayer itself was not correct - rather than asking the Lord for assistance I questioned Him.  I resolved myself to overcome my affliction, I began running; working myself up to a 6-1/2 minute mile.  I participated in mini-marathons, training marathon distance, weight training.  Overtime, through fortitude and determination I got stronger.  I convinced myself that the Lord was an observer and not a participant in my life.  I drifted from the Church and came to despise any weakness in myself - weakness was something to remedied at all costs.

I now realize the purpose of the first forge was to develop the determination and will that would enable me to withstand the fire of the second forge.  I slowly returned to Church, prayed my Rosary daily and I began to understand:

The adversity we face has a purpose.  The Bible does not frame adversity and suffering as an accident or a mistake.  It treats it as a deliberate refining process.  The prophet Malachi describes the Lord as a "refiner's fire":
 
Malachi 3:2-3 is not just poetic language.  Rather, it reflects the mechanics of how God forms a man.   The fire does not destroy the silver -- it burns away the dross so that what remains is pure.  The furnace does not consume the gold -- it separates the impurities so the metal can be shaped and used.

The same principle appears again and again throughout Scripture. Job was stripped of everything so his faith could be proven genuine. David cried out in the Valley of the Shadow and emerged declaring God’s name in the great assembly. Peter wrote that these trials come “ so that the proven genuineness of your faith — of greater worth than gold, which perishes even though refined by fire — may result in praise, glory and honor when Jesus Christ is revealed ” (1 Peter 1:7).  The Forge is not random suffering. It is the deliberate, holy fire of the Refiner.

The Lord Himself sits as a refiner and purifier of silver. He does not watch from a distance — He is in the fire with us:

“But who can endure the day of his coming? Who can stand when he appears? For he will be like a refiner’s fire or a launderer’s soap. He will sit as a refiner and purifier of silver.” (Malachi 3:2–3)

This is the ancient pattern repeated throughout Scripture. Isaiah declares, “ See, I have refined you, though not as silver; I have tested you in the furnace of affliction ” (Isaiah 48:10). 

Zechariah promises the Lord will put His people through the fire “ and refine them as silver is refined, and test them as gold is tested ” (Zechariah 13:9). Peter tells us these trials come so that “ the proven genuineness of your faith—of greater worth than gold, which perishes even though refined by fire—may result in praise, glory and honor when Jesus Christ is revealed ” (1 Peter 1:6–7).
The Forge is not punishment. 

It is purification. The enemy attacks what is most important — first a man's spirit and then his physical strength — thinking he will break the man. But every blow he lands only turns up the heat of the Forge the Lord Himself placed us in. The small god of this world believes he is destroying us. The King knows he is only intensifying the refining process. More dross is burned away and the steel that emerges is unbreakable.

Nothing the enemy takes is beyond the Lord’s reach.

Job, stripped of everything, declared, “But he knows the way that I take; when he has tested me, I will come forth as gold” (Job 23:10). David walked through the Valley of the Shadow and emerged declaring God’s name in the great assembly. The same fire that tested them is the fire that tested me.

This is the biblical reality. This is what was happening on July 5, 2022, when the disease had ravaged my body and the voice spoke clearly into the darkness:

“Bench. Squat. Deadlift.”

The Valley had become the Forge and the Forge and faith are not separate. They are locked in a living covenant. The fire tests the faith, and the faith steps willingly into the fire. Each strengthens the other.  The Lord does not call us to passive endurance.  He calls us to active obedience inside the furnace.  This is why the voice in the recliner did not say “pray harder” or “wait for healing.” It said something far more demanding.

That command was the birth of the Iron Church — the physical expression of faith when everything in the body is screaming to quit. The Forge is the fire.  The Iron Church is the hammer and the anvil.  Together they shape the man.  

Faith (Jesus's Church) without Iron Church becomes soft and theoretical,  Iron Church without Jesus's Church becomes mere exercise and pride.  But, when the two are yoked together, iron sharpens iron (Proverbs 27:17).   Each strengthens the other, and something unbreakable is formed and men rise to that for which they were made -- dominion over they physical world.

Once one begins to understand this, every rep, every set, every mile on the treadmill becomes an act of worship.  Each time the body wants to stop and the will chooses to continue, faith is proven genuine and the dross is burned away.  The enemy thought he could destroy me by attacking the body.  Instead, I was driven deeper into the Forge and forced me to build the very thing that would make me stronger — deeper faith God Church yoked to the Iron Church.

This is the mystery of the refining fire: the hotter the flame, the stronger the steel. The greater the adversity, the deeper the faith. The two do not oppose each other. They complete each other.

And so the Valley became the Forge, and the Forge called forth the Iron Church.




\chapter{The Valley as Forge – The Same Fire}
[Psalm 23:4]

\chapter{The Locusts and the Repayment}
[Joel 2:25]

\chapter{Betrayal in the Valley}
[Psalm 41:9]

\chapter{The Self-Built Prison}
[Proverbs 14:1, 2 Timothy 3:1–5]

\chapter{Walking Through Without Bitterness}
[Job’s perseverance, Ephesians 4:31–32]

\chapter{The Three Pillars of Restoration}

\chapter{The Final Vindication – Repaid Double}
[Job 42:10–17]

% Conclusion
\chapter*{Conclusion}
\addcontentsline{toc}{chapter}{Conclusion}

Revelation 21:5 – “Behold, I make all things new.”

The locusts came.  
I cried out.  
God answered.  
The repayment is unfolding — one rep, one breath, one day at a time.

\end{document}

